\section{Implementation Details for Metadata and Curation} 
\label{appx:details}

\subsection{Unigram and Bigram Tokenizer for Special Languages}
\label{appx:tokenizers}

Most modern languages around the world adopt writing systems that use ``spaces'' to separate words, except for some of the Asian languages, known as ``scriptio continua''\citep{scriptio_continua}. We find several open source tokenizers for many of these languages developed by local communities, as shown in the Table~\ref{tbl:tokenizer}, in order to properly split text into words while preserve semantic integrity. Note these tokenizers are only used to process Wikipedia dump labeled with the listed wiki codes (e.g., not on alt-texts' substring matching).

\begin{table}[h]
\centering
\scalebox{0.85}{
\begin{tabular}{c|c}
\toprule
\textbf{Wiki Code} & \textbf{Tokenizer Name} \\
\midrule
bo,dz & Tibetan Tokenizer\\
ja,ryu  & Japanese Tokenizer\\
km  & Khmer Tokenizer\\
lo & Lao Tokenizer\\
my & Myanmar Tokenizer\\
th & Thai Tokenizer\\
zh,zh\_classical,zh\_yue & Chinese Tokenizer\\
\bottomrule
\end{tabular}
}
\vspace{1.5mm}
\caption{Tokenizers for special languages.}
\label{tbl:tokenizer}
\end{table}

\subsection{Scaling Curation} 
\label{sec:impl_curation}
Worldwide scaling of data curation significantly increases time and space complexity due to store metadata across hundreds of languages. To efficiently handle this complexity, we leverage several efficient implementation to enable worldwide curation (that English only curation does not have):

\begin{itemize}
    \item \textbf{Efficient String Matching}: We adopt the Aho-Corasick algorithm~\footnote{\href{https://en.wikipedia.org/wiki/Aho-Corasick_algorithm}{\tt https://en.wikipedia.org/wiki/Aho-Corasick\_algorithm}}\footnote{\url{https://pypi.org/project/pyahocorasick}}, which utilizes prefix trees (tries), for rapid substring matching. The matching speed is about 2k times faster than Meta CLIP's brute-force implementation, enabling matching with million-scale metadata.
    \item \textbf{Lazy Metadata Loading}: We pre-build and store the metadata into an Aho-Corasick automaton for each language separately, loading these automaton dynamically and only when encountering a new language for alt-text during processing, thereby minimizing the total number of languages encountered for each shard of data and saving re-compiling time for automation on a new shard.
    \item \textbf{Memory Management for Probabilities}: To address memory constraints during sampling for balancing, we utilize memory-mapped file loading (mmap) to efficiently access counts per entry across all languages, preventing out-of-memory errors caused by loading all the counts from different languages.
\end{itemize}

These implementation choices ensure the worldwide data curation is computationally feasible and scalable to billions of image-text pairs from hundreds of languages. 

\paragraph{Mitigation and Benchmark Deduplication} We run a state-of-the-art safety classifier to remove NSFW contents (e.g., adult, sexual, violence) from training data. We also apply face detector to remove human biometrics and personally identifiable information from data. To avoid benchmark leakage, we remove any overlap with ImageNet evaluation sets by performing deduplication using 64-bit hashes. These hashes are generated by applying random projection to feature embeddings from a similarity search model, reducing them to 64 dimensions followed by sign-based quantization.

\section{Training Setup}
\label{sec:training_setup}
To remove confounding factors and generalize our findings, we follow OpenAI CLIP and Meta CLIP training setup with changes for worldwide scaling, detailed in Table~\ref{tbl:hp}.  
Our data curation algorithm is running in parallel with 800 jobs (each job has 40GB CPU memory) and it takes 1 hour to substring match and count for all alt-text pairs.

\begin{table}[h]
\centering
\scalebox{0.8}{
    \setlength\tabcolsep{3.0pt}
    \begin{tabular}{l| c | c}
    \toprule 
    Hyperparameter & OpenAI CLIP / Meta CLIP & Meta CLIP~2\\ 
    \midrule
        Activation Function & QuickGELU & QuickGELU \\
        Seen Pairs & 12.8B & 29B (2.3$\times$) \\
        Batch Size & 32768 & 75366 (2.3$\times$) \\
        Learning Rate & 4.0e-4 (L/14, H/14) & 4.0e-4 (H/14) \\
        Warm-up & 2k & 2k \\
    \bottomrule
    \end{tabular}
}

\caption{Hyperparameters of OpenAI CLIP / Meta CLIP vs Meta CLIP~2.}
\label{tbl:hp}
\end{table}

\section{Limitation on Benchmark}
\label{sec:lim_bench}

High-quality benchmarks are essential for researchers to understand the efficacy of proposed changes. After decades of meticulous efforts, the community has established reliable and diverse datasets to enable research advancement in vision and multimodal areas~\citep{deng2009imagenet,russakovsky2015imagenet,radford2021learning}. However, these datasets consist mainly of content scraped from North America and Western Europe (NA+EU) and focus on English~\citep{shankar2017no,de2019does}. It is a long and resource-intensive endeavor to build similar benchmarks for unbiased and comprehensive evaluation of worldwide data and resulting representations, for the world outside NA+EU or English-speaking community, due to the complexity of covering diverse concepts across geo-locations, cultures, and languages. XM3600~\citep{thapliyal2022crossmodal} aims to build geographically diverse datasets by selecting images from Open Images Dataset~\citep{kuznetsova2020open} based on metadata of GPS coordinates, but later research~\citep{pouget2024no} suggests Open Images Dataset is biased towards Western images or specific activities (e.g., tourism). 
GeoDE~\citep{ramaswamy2023geode} recruits human workers on crowdsourcing platform to collect geographically diverse images for predefined object classes. Crowdsourcing is an economic way to collect human annotations, but the demographic background and proficiency of the workers are not guaranteed, nor is the quality of the collected data.
Few efforts such as CVQA~\citep{mogrovejo2024cvqa}
attempt to scale annotation and control quality simultaneously 
by utilizing experts in machine learning community or existing materials as seeds. These efforts offer relatively unbiased evaluation with reasonable coverage in capabilities (e.g., cultural diversity, multimodal problem solving for exam questions across countries) of interests. We believe benchmarks of similar quality but built for evaluating more general and comprehensive capabilities will reveal the true potential of worldwide data and resulting representations developed in this work.
